% Forest of Knowledge — Whitepaper (Bitcoin-style)
% Companion to the full technical paper (mogae-paper-v5.tex)

\documentclass[11pt]{article}
\usepackage[margin=1in]{geometry}
\usepackage[utf8]{inputenc}
\usepackage[T1]{fontenc}
\usepackage{lmodern}
\usepackage{hyperref}
\usepackage{booktabs}
\usepackage{enumitem}
\usepackage{amsmath}
\usepackage{graphicx}

\title{\textbf{Forest of Knowledge}\\[4pt]
\large Separating What a Model Knows from How It Thinks}

\author{Erik de Bruijn\\
\texttt{erik@erikdebruijn.nl}}

\date{March 2026}

\begin{document}
\maketitle

% ============================================================================
\begin{abstract}
Current language models entangle reasoning ability with factual knowledge in
one monolithic block of parameters. Adding knowledge means adding parameters,
which means more memory, more energy, and more centralized control. We
propose \textbf{Forest of Knowledge}: an architecture that separates
\emph{fluid intelligence} (reasoning, syntax, composition---shared by all
tasks) from \emph{crystallized intelligence} (domain-specific facts,
terminology, patterns---different per specialty) using lightweight,
hot-pluggable adapters measured in kilobytes. The fluid trunk stays fixed.
Specialists are trained independently, stored on cheap flash, loaded on
demand, and contributed by a distributed community. Across 50+ experiments
on Qwen3-8B, we show that this separation works: 8 experts grown
autonomously via a difficulty signal (\emph{learntropy}) reduce perplexity
by 22.8\% while adding less than 1\% to model size. The architecture enables
a paradigm shift from centralized model training to distributed,
community-driven specialization.
\end{abstract}

% ============================================================================
\section{The Problem}

Language models today are monoliths. A single set of parameters encodes
everything: grammar, reasoning, medical terminology, legal precedent,
programming syntax, and conversational slang. Training such a model requires
thousands of GPUs, months of compute, and centralized control by a handful
of frontier labs. The barrier to entry is not technical insight but capital:
training a single frontier model costs tens of millions of dollars. This
concentrates the power to decide \emph{what knowledge a model contains} in
the hands of a few organizations.

This creates four problems:
\begin{enumerate}[nosep]
\item \textbf{Adding knowledge is expensive---for training and inference.}
      Every new specialty requires retraining or fine-tuning the entire
      model, and the resulting larger model is more expensive to \emph{run}
      for every query, not just to train. A hospital that wants medical
      expertise must either pay for a larger, slower model or accept a
      general one.
\item \textbf{Memory is the bottleneck.} The DRAM needed to hold these
      models has more than doubled in price over the past year. A 70B model
      requires 140\,GB of memory at half precision---more than most hardware
      can provide.
\item \textbf{Knowledge is centralized.} Only organizations with massive
      compute budgets can train frontier models. The rest of the world
      consumes what those labs choose to produce. Domain experts---doctors,
      lawyers, engineers---have no mechanism to contribute their expertise
      directly to the models they use.
\item \textbf{No incentive for contribution.} Even if distributed training
      were possible, there is no mechanism to reward contributors for their
      compute and expertise, verify the quality of their contributions, or
      protect the system against malicious inputs.
\end{enumerate}

% ============================================================================
\section{The Idea}

We observe that language models contain two fundamentally different kinds of
capability:

\begin{itemize}[nosep]
\item \textbf{Fluid intelligence}: reasoning, logical composition,
      syntactic structure, attention patterns. This is shared across all
      domains and changes slowly. It lives in the model's ``trunk''---the
      shared layers that every input passes through.
\item \textbf{Crystallized intelligence}: medical facts, legal definitions,
      programming idioms, cultural references. This is domain-specific,
      changes frequently, and can be enormous in aggregate.
\end{itemize}

Our key insight: \emph{these don't need to live in the same parameters.}

The trunk (fluid intelligence) stays fixed after initial training. Domain
knowledge lives in tiny LoRA adapters---lightweight modifications to the
trunk's feed-forward layers that encode specialized corrections. A rank-4
adapter is 4\,KB. A rank-64 adapter is 256\,KB. Even thousands of
specialists total less than a gigabyte.

% ============================================================================
\section{How It Works}

\subsection{The Trunk}

We start with a strong pretrained model (e.g., Qwen3-8B). This model
already possesses fluid intelligence: it can reason, compose sentences, and
follow instructions. We freeze its parameters permanently. The trunk is the
shared root system that all specialists inherit.

\subsection{Growing Experts}

Specialists grow via a process inspired by Piaget's theory of cognitive
development:

\begin{enumerate}[nosep]
\item \textbf{Start small.} Begin with one adapter at low rank (4
      dimensions out of 4096). This adapter learns to make small corrections
      to the trunk's output.
\item \textbf{Measure difficulty.} For each token, measure how hard it is
      for the current expert (\emph{learntropy}---per-token cross-entropy
      weighted by routing). High learntropy means the expert is struggling.
\item \textbf{Split when struggling.} When an expert's difficulty stays
      high, split it into two: one inherits the parent's stable knowledge,
      the other explores the hard tokens.
\item \textbf{Increase rank when needed.} If difficulty persists after
      splitting, the expert needs more capacity. Double its rank
      (4$\to$8$\to$16$\to$64). Our experiments show the system converges
      to rank-64 autonomously---exactly the empirically optimal point.
\item \textbf{Repeat.} The tree grows from 2 experts to 8 over 20K training
      steps, with all 8 active and contributing. Reproduced across seeds.
\end{enumerate}

\subsection{Routing}

A lightweight Gumbel-softmax router decides which expert(s) to activate for
each token. At inference time, routing is nearly instant---a single linear
layer per position. Only activated experts consume compute.

% ============================================================================
\section{From Centralized Training to Distributed Specialization}

The current paradigm mirrors the pre-Bitcoin financial system: a small
number of centralized entities must train ever-larger monolithic models.

Forest of Knowledge proposes the alternative:

\begin{center}
\begin{tabular}{@{}lll@{}}
\toprule
& \textbf{Centralized (current)} & \textbf{Forest of Knowledge} \\
\midrule
Training & One lab, thousands of GPUs & Community, commodity hardware \\
Knowledge & Baked into parameters & Hot-pluggable adapters \\
Adding expertise & Retrain the model & Train one adapter (KB-scale) \\
Memory & All in expensive DRAM & Trunk in DRAM, rest on flash \\
Cost scaling & Linear in model size & Sub-linear (adapters are tiny) \\
Governance & Lab decides what to learn & Community decides \\
\bottomrule
\end{tabular}
\end{center}

Each adapter trains independently. A contributor needs only the frozen trunk
(a one-time download) and a few kilobytes of adapter weights. They train
their specialist on domain-specific data using learntropy signals, then
return the updated adapter. No synchronization between contributors is
needed.

This is analogous to Bitcoin mining: a massively parallel, compute-bound
effort where independent participants contribute to a shared system. The
difference is that our ``proof of work'' is \emph{proof of useful learning}:
every adapter that reduces perplexity on its domain adds genuine value.

% ============================================================================
\section{Verification and Incentives}

Distributed training requires mechanisms to verify contributions and
incentivize honest participation---the same challenge Bitcoin solves with
proof-of-work and block rewards.

\paragraph{Proof of useful learning.}
Unlike Bitcoin's proof-of-work (which produces no useful output beyond
consensus), adapter training produces a measurable artifact: an adapter that
must demonstrably reduce perplexity on a held-out validation set for its
target domain. Verification is asymmetric: training an adapter costs hours
of GPU compute, but validating it costs seconds (evaluate PPL on a standard
benchmark). This asymmetry makes fraud expensive and detection cheap.

\paragraph{Data integrity.}
Training data is pre-agreed: specific shards of curated datasets (e.g.,
FineWeb-2 medical subset). SHA-256 checksums of training batches are
published alongside the adapter. Any validator can confirm the adapter was
trained on the claimed data by re-running a small number of spot-check
steps with the same seed and verifying intermediate checkpoints match.

\paragraph{Malicious adapter detection.}
Poisoned adapters (backdoors, biased outputs) are detected via standard
federated learning techniques: outlier detection on weight norms, gradient
clipping bounds, and behavioral testing on adversarial inputs. Because
adapters are kilobyte-scale, exhaustive behavioral auditing is feasible.

\paragraph{Incentive mechanism.}
The forest is a public good: all adapters are open weights. The primary
incentive is \emph{self-interest}: a cardiologist who trains a cardiology
adapter benefits directly from a better model for their own work. A lawyer
who trains a contract-law adapter gets a tool that understands their
domain. The contribution makes the entire forest stronger as a
side-effect---just as Wikipedia grows because contributors want articles
about topics they care about, not because they are paid.

This market-driven approach is sufficient for an MVP. As the ecosystem
matures, additional mechanisms can emerge: token rewards to compensate
compute costs for contributors who train adapters outside their own domain,
reputation systems that give active contributors governance weight over
trunk updates, and backwards-compatibility guarantees for high-value
adapters (once your adapter is widely used, trunk updates must remain
compatible with it---making your contribution a protected stake in the
system).

% ============================================================================
\section{Why It's Cheaper}

Three properties make this dramatically cheaper than monolithic models:

\paragraph{Storage economics.}
SSD/flash storage is $\sim$80$\times$ cheaper per GB than DRAM. HDDs are
$>$100$\times$ cheaper. A forest of 10,000 rank-64 specialists requires
2.5\,GB of flash---less than a single smartphone app. The trunk (the only
part requiring DRAM) stays small.

\paragraph{Energy proportionality.}
Only activated adapters consume compute. If a query is about medicine, only
the medical specialist is loaded. Legal, code, and conversational experts
remain on flash, consuming zero energy. Training new specialists requires
only adapter-scale compute, not full-model retraining.

\paragraph{DRAM crisis mitigation.}
DRAM prices have more than doubled in the past year. Forest of Knowledge
sidesteps this: the bulk of knowledge lives on cheap flash, not expensive
DRAM. The trunk occupies a fixed DRAM footprint regardless of how many
specialists exist.

% ============================================================================
\section{Key Results}

Across 50+ experiments on Qwen3-8B (see technical paper for full details):

\begin{enumerate}[nosep]
\item \textbf{8 experts, all active.} Learntropy-directed splitting grows
      2$\to$8 experts with all 8 contributing. Reproduced across seeds
      ($M_{ij}$ CV up to 2.34).
\item \textbf{22.8\% perplexity reduction} compared to the base model,
      with adapters totaling $<$1\% of model parameters.
\item \textbf{Autonomous rank discovery.} The system converges from rank-4
      to rank-64 through learntropy-driven doublings---finding the empirical
      optimum without manual tuning.
\item \textbf{Speculative splitting with rollback.} Instead of
      predicting when to split, the system speculatively splits every
      $N$ steps and reverts if loss doesn't improve. In 25K steps, only
      2 of 8 attempted splits were kept (at steps 9K and 21K)---the
      system autonomously determined that 4 experts suffice.
\item \textbf{Hot-plug demo.} A Wing Chun martial arts adapter (6\,MB)
      loads in 39\,ms. The base model gives generic, repetitive answers;
      the adapter produces specific knowledge including Cantonese
      terminology. Unloading restores the base model identically.
\item \textbf{Domain selectivity peaks at rank-64} (1.6\% of hidden
      dimension). Lower ranks specialize well but have limited capacity;
      higher ranks lose specialization pressure.
\item \textbf{Zipf-distributed importance.} Expert ablation damage follows
      a power law: one backbone expert is critical ($\Delta$PPL +43), a few
      mid-tier experts are important (+3 to +17), and the long tail
      contributes incrementally (+0.1 to +3). This matches the natural
      access pattern for tiered storage.
\end{enumerate}

% ============================================================================
\section{``I Don't Know, Yet'': Hallucination as a Routing Problem}

Current language models hallucinate because they have no mechanism to
distinguish ``I know this'' from ``I'm guessing.'' The model always
generates a confident answer, whether or not the knowledge exists in its
parameters.

Forest of Knowledge reframes hallucination as a \emph{routing problem}.
Learntropy---the per-token difficulty signal that drives all adaptation
decisions---also serves as a real-time uncertainty indicator:

\begin{itemize}[nosep]
\item \textbf{Low learntropy} on a query: the trunk or an active adapter
      handles it well. Generate normally.
\item \textbf{High learntropy, adapter available}: the model lacks the
      knowledge but a verified adapter exists. Response: ``I don't have
      this knowledge loaded. Found: {$\surd$} \emph{Martial Arts:
      Wing Chun} (64\,KB). Load it?''
\item \textbf{High learntropy, no adapter}: the knowledge doesn't exist
      in the forest yet. Response: ``I don't know. No verified knowledge
      module exists for this topic yet.''
\end{itemize}

This transforms the hallucination problem from an unsolvable alignment
challenge into a measurable, mechanistic decision: \emph{if learntropy is
high and no adapter reduces it, stop generating and say so.} The model
can even proactively discover and prefetch relevant adapters from a
registry by matching the query embedding against adapter capability
signatures---loading knowledge on demand, in milliseconds, without ever
hallucinating.

The long tail of crystallized knowledge that causes most hallucinations
(rare facts, niche terminology, recent events) is precisely what the
adapter forest is designed to store. As the forest grows through community
contributions, the space where the model must say ``I don't know'' shrinks
continuously---without retraining the trunk.

\emph{Experimental progress:} Initial experiments show that the per-layer
delta gate itself acts as an ``I don't know'' detector. When a domain adapter
encounters topics it was not trained on, the gate drops to near-generic levels
(2 of 3 adapters tested show known/unknown ratios of 2.25--4.11$\times$;
the third shows a weaker 1.86$\times$ signal, likely due to domain overlap
with generic text). This is a learned behavior, not an architectural
add-on---the gate discriminates because it was trained to activate selectively.

\paragraph{Demand-driven growth.}
The ``I don't know'' signal has a second use: it reveals where the grove
has gaps. When many queries trigger low gate activation across all adapters,
this is a \emph{cache miss}---a topic that users need but no adapter covers.
Logging these cache misses creates a demand signal for training priorities:
the grove grows where users need it, not where contributors guess. This is
analogous to how CDNs cache popular content---except here the ``cache'' is
a trained adapter and the ``miss'' triggers training, not just a round-trip
to origin.

\emph{Caveat:} The gate-based IDK detection is experimentally supported
(3 adapters, multiple unknown domains) but not yet validated as a reliable
hallucination preventer. Hallucination is harder than uncertainty: models
can be confidently wrong. Whether the gate signal correlates with factual
correctness---rather than mere domain membership---requires further testing.
The demand-driven pipeline is a theoretical extension, not yet validated
end-to-end.

% ============================================================================
\section{Implications}

\paragraph{For the planet.}
Training adapters instead of full models reduces energy consumption by
orders of magnitude. A single adapter trains in minutes on one GPU;
a full model trains in months on thousands. Storing knowledge on flash
instead of DRAM reduces hardware costs and e-waste. The efficiency gains
compound: a forest of 10,000 specialists is simultaneously more
knowledgeable and cheaper to run than a single monolithic model trained
to cover the same domains.

\paragraph{For the community.}
Anyone with a GPU and domain expertise can train a specialist. The forest
grows through contribution, not centralization. A medical researcher trains
a cardiology adapter. A lawyer trains a contract-law adapter. A programmer
trains a Rust adapter. Each returns a kilobyte-scale module that plugs into
the shared trunk.

\paragraph{For practitioners.}
Hot-pluggable adapters mean you can add medical, legal, or domain-specific
expertise to any compatible trunk without retraining. Adapters are measured
in kilobytes and load in microseconds from NVMe. Inference cost depends
only on the trunk size plus the few active adapters---not on the total
knowledge in the forest.

\paragraph{For researchers.}
The separation of fluid and crystallized intelligence provides a clean
experimental framework. You can study reasoning (trunk) and knowledge
(adapters) independently, measure specialization via ablation, and grow
expert trees that reveal the structure of what models learn.

% ============================================================================
\section{Conclusion}

Language models don't need to be monoliths. Fluid intelligence (reasoning)
and crystallized intelligence (knowledge) can be separated, with knowledge
living in tiny, hot-pluggable adapters that a distributed community trains
independently. The learntropy signal---a measure of per-token
difficulty---drives all adaptation decisions: when to split, where to
specialize, how fast to learn, and when to increase capacity.

The result is a forest: a shared root system of fluid intelligence,
supporting an ever-growing canopy of crystallized knowledge. Each tree in
the forest is independently trainable, independently loadable, and
independently valuable. The architecture makes knowledge addition
\emph{cheap}, \emph{distributed}, and \emph{composable}---a paradigm shift
from centralized training to collaborative specialization.

\vspace{12pt}
\noindent\emph{Full technical details, experimental methodology, and all
50+ experiment results are available in the companion technical paper:
``Forest of Knowledge: Learntropy-Driven Separation of Fluid and
Crystallized Intelligence in Language Models'' (de Bruijn, 2026).}

\end{document}
